% preamble.tex

\usepackage[utf8]{inputenc}
\usepackage[T1]{fontenc}
\usepackage[english]{babel} % Ho cambiato in italian, supponendo che la maggior parte del testo sia in italiano
\usepackage{amsmath}
\usepackage{amssymb}
\usepackage{graphicx}
\usepackage{hyperref}
\usepackage{geometry}
\usepackage{fontspec} % Utile per Xetex/LuaLaTeX ma mantenuto
\usepackage{booktabs}
\usepackage{amsthm}
\usepackage{enumitem}
\usepackage{caption}
\usepackage{thmtools}
\usepackage{standalone} % Già presente nel tuo elenco
\usepackage{tkz-euclide}
\usepackage{graphicx}
\usepackage[T1]{fontenc}
\usepackage{wesa}

% --- PACCHETTI AGGIUNTI PER LA FORMATTAZIONE ---
\usepackage{parskip}  % Aggiunge spazio tra i paragrafi (stile "blocco")
\usepackage{fancyhdr} % Per intestazioni e piè di pagina personalizzati

% --- IMPOSTAZIONI GEOMETRIA E FONT ---
\geometry{a4paper, margin=1in}

% Configura lo stile 'plain' (usato da \maketitle e \tableofcontents)
% per essere coerente con il piè di pagina
\fancypagestyle{plain}{
    \fancyhf{} % Pulisce
    \cfoot{\thepage} % Mette solo il numero di pagina
    \renewcommand{\headrulewidth}{0pt} % Niente linea su queste pagine
}

% Configurazione Header/Footer per le altre pagine
\pagestyle{fancy}
\fancyhead{} % Pulisce l'intestazione
\fancyfoot{} % Pulisce il piè di pagina
\fancyhead[L]{\nouppercase{\rightmark}} % Nome Sezione a sinistra (o Chapter)
\fancyhead[R]{\nouppercase{\leftmark}}  % Nome Capitolo a destra (o Sezione)
\fancyfoot[C]{\thepage} % Numero di pagina al centro
\renewcommand{\headrulewidth}{0.4pt}
\renewcommand{\footrulewidth}{0pt}


% --- DEFINIZIONI TEOREMI ---
% (Questa parte è cruciale per la tua struttura)
\newtheoremstyle{mythmstyle}%
    {}%
    {}%
    {\itshape}%
    {}%
    {\bfseries}%
    {}%
    { }%
    {\thmname{#1}\thmnumber{ #2}%
    \thmnote{: #3\addcontentsline{toc}{subsubsection}{{\itshape#1}: #3}}. }

% 2. Imposta questo stile come quello ATTUALE
\theoremstyle{mythmstyle}

\newtheorem{theorem}{Th}[section]
\newtheorem{definition}{Def}[section]
\newtheorem{proposition}{Prop}[section]
\newtheorem{corollary}{Cor}[section]
\newtheorem{algorithm}{Alg}[section]
\newtheorem{lemma}{Lemma}[section]
\newtheorem*{remark}{Rem}

% Definizioni utili per il Main file
\newcommand{\inputlesson}[1]{\clearpage\input{#1}\clearpage} % Comando per input lezione